\chapter{Bentuk Hermitian}\label{ch:b2:hermitian}

Ruang vektor kompleks punya geometri hasil kali dalamnya sendiri, dengan
satu putaran: linear pada satu peubah, dan linear-\emph{sekawan} pada
peubah yang lain. Imbalan karena menerima putaran itu adalah teori spektral
yang bahkan lebih bersih daripada yang real --- sebab \omterm{def:b2:hermitian:adjoint}{endomorfisma Hermitian}
bernilai eigen real, yang \omterm{def:b2:hermitian:adjoint}{uniter} bernilai eigen bermodulus satu, dan keduanya
\omterm{def:b2:reduction:diag}{dapat didiagonalkan} dalam basis ortonormal. Bab pendek ini menjalankan
program Euklides \cref{ch:b2:quadratic} atas $\C$.

\section{Hasil kali dalam Hermitian}

\begin{definition}\label{def:b2:hermitian:def}
Sebuah \emph{hasil kali dalam Hermitian}\index{hasil kali dalam Hermitian}
pada ruang vektor kompleks $E$ adalah pemetaan $\langle\cdot,\cdot\rangle
\colon E \times E \to \C$ yang linear pada peubah keduanya,
simetrik-sekawan (yakni $\langle y, x\rangle = \conj{\langle x,
y\rangle}$ --- sehingga linear-sekawan pada peubah pertamanya), dan
definit positif (yakni $\langle x, x\rangle > 0$ untuk $x \neq 0$). Adapun
contoh bakunya pada $\C^n$:
\[
\langle x, y \rangle = \sum_{i=1}^{n} \conj{x_i}\, y_i ;
\]
sedangkan pada fungsi \omterm{def:b2:metric:continuity}{kontinu}, $\langle f, g\rangle = \int_a^b \conj
f\,g$. \omterm{def:b2:nvs:norm}{Normanya}: $\norm x = \sqrt{\langle x,x\rangle}$; dan ruang
kompleks berdimensi hingga yang dilengkapi demikian disebut \emph{ruang
Hermitian}\index{ruang Hermitian}.
\end{definition}

\begin{example}[Perhitungan pertama]\label{ex:b2:hermitian:firstcomp}
Di $\C^2$, ambillah $x = (1+\iu,\ 2-\iu)$ dan $y = (\iu,\ 1)$.
Maka
\[
\norm x^2 = \abs{1+\iu}^2 + \abs{2-\iu}^2 = 2 + 5 = 7,
\qquad
\norm y^2 = 1 + 1 = 2,
\]
lalu, dengan mengawankan argumen pertamanya,
\[
\langle x, y\rangle
= \conj{(1+\iu)}\,\iu + \conj{(2-\iu)}\cdot1
= (1-\iu)\iu + (2+\iu)
= (\iu + 1) + (2 + \iu) = 3 + 2\iu .
\]
Cauchy--Schwarz terperiksa: $\abs{\langle x, y\rangle}^2 = 9 +
4 = 13 \leq 14 = \norm x^2\norm y^2$ --- yang dekat ke kesamaan,
karena $x$ dekat ke kelipatan $y$. Perhatikan pula
$\langle y, x\rangle = \conj{3 + 2\iu} = 3 - 2\iu$: yakni
kesimetrikan-sekawan yang beraksi, dan itulah sebabnya $\langle x,
x\rangle$ selalu real.
\end{example}

\begin{theorem}[Cauchy--Schwarz, versi kompleks]\label{thm:b2:hermitian:cs}
$\abs{\langle x, y\rangle} \leq \norm x \norm y$, dengan kesamaan bila dan
hanya bila $x, y$ bergantung linear; dan $\norm\cdot$ merupakan \omterm{def:b2:nvs:norm}{norma}. Lebih
jauh basis ortonormal ada (sebab Gram--Schmidt berjalan kata demi kata), dengan
\[
x = \sum_i \langle e_i, x\rangle\, e_i,
\qquad
\norm x^2 = \sum_i \abs{\langle e_i, x\rangle}^2 .
\]
\end{theorem}

\begin{proof}
Untuk $y \neq 0$ dan $t \in \C$: $0 \leq \norm{x - ty}^2 = \norm x^2
- 2\Re\bigl(\conj t\langle y, x\rangle\bigr) + \abs t^2\norm y^2$.
Pilihlah $t = \frac{\langle y, x\rangle}{\norm y^2}$:
\[
0 \leq \norm x^2 - \frac{\abs{\langle y, x\rangle}^2}{\norm y^2},
\]
yang tak lain ketaksamaannya; sedangkan kesamaannya memaksa $x = ty$.
Ketaksamaan segitiganya, selengkapnya:
\[
\norm{x + y}^2 = \norm x^2 + 2\,\Re\langle x, y\rangle
+ \norm y^2
\leq \norm x^2 + 2\,\abs{\langle x, y\rangle} + \norm y^2
\leq \bigl(\norm x + \norm y\bigr)^2 ,
\]
dengan memakai $\Re z \leq \abs z$ lalu Cauchy--Schwarz; adapun kehomogenan
dan pemisahannya seketika, jadi $\norm\cdot$ merupakan \omterm{def:b2:nvs:norm}{norma}.
Gram--Schmidt: seperti pada kasus realnya, dengan sekawannya
ditempatkan menurut definisinya (perhatikan konvensinya: hasil kali kita
bersifat linear-sekawan pada slot \emph{pertama}, jadi koordinatnya adalah
$\langle e_i, x\rangle$; dan contoh berikutnya menjalankan algoritmanya
sekali secara \omterm{def:b2:metric:complete}{lengkap}).
\end{proof}

\begin{example}[Gram--Schmidt kompleks, dijalankan lengkap]\label{ex:b2:hermitian:gramschmidt}
Ortonormalkan basis $v_1 = (1, \iu)$, $v_2 = (0, 1)$ milik
$\C^2$. Vektor pertamanya: $\norm{v_1}^2 = \abs1^2 + \abs\iu^2 = 2$,
jadi $e_1 = \frac{1}{\sqrt2}(1, \iu)$. Proyeksikan $v_2$ --- dengan
sekawannya pada slot pertama:
\[
\langle e_1, v_2\rangle
= \frac{1}{\sqrt2}\bigl(\conj{1}\cdot0 +
\conj{\iu}\cdot1\bigr)
= \frac{-\iu}{\sqrt2} ,
\qquad
v_2 - \langle e_1, v_2\rangle e_1
= (0,1) + \frac{\iu}{2}\,(1, \iu)
= \Bigl(\frac\iu2,\ \frac12\Bigr) .
\]
\omterm{def:b2:nvs:norm}{Normanya} $\sqrt{\frac14 + \frac14} = \frac{1}{\sqrt2}$, jadi
$e_2 = \frac{1}{\sqrt2}(\iu, 1)$. Periksa: $\langle e_1,
e_2\rangle = \frac12(\conj1\cdot\iu + \conj\iu\cdot1) =
\frac12(\iu - \iu) = 0$. Koordinat $v_2$ dalam basis barunya:
$v_2 = \langle e_1, v_2\rangle e_1 + \langle e_2, v_2\rangle
e_2$ dengan $\langle e_2, v_2\rangle = \frac{1}{\sqrt2}$ ---
perhatikan urutannya: sebab $\langle v_2, e_1\rangle$ akan memberi
koefisien \emph{sekawannya}. Pelajaran penutupnya: algoritmanya adalah
algoritma Euklides kata demi kata; satu-satunya jebakannya adalah tempat
jatuhnya pengawanan itu, dan menghitung $\norm{v_2 -
\text{proy}}^2 > 0$ diam-diam memakai kepositifannya --- yakni aksioma
yang membuat seluruh geometrinya bekerja.
\end{example}

\section{Adjoin, endomorfisma Hermitian dan uniter}

\begin{definition}\label{def:b2:hermitian:adjoint}
Adapun \emph{adjoin} $u^*$ milik $u \in \mathcal{L}(E)$ didefinisikan lewat
$\langle u^*(x), y\rangle = \langle x, u(y)\rangle$; dan dalam basis
ortonormal berlaku $\operatorname{Mat}(u^*) = \conj{A}^{\mathsf T} =:
A^{\dagger}$ (yakni \omterm{def:b2:linalg:transpose}{transpos} sekawannya) --- sebab, bila $B =
(b_{ij})$ matriks $u^*$ dalam basis ortonormal
$(e_i)$, maka $b_{ij} = \langle e_i, u^*(e_j)\rangle$, lalu
kesamaan pendefinisinya memberi
\[
\conj{b_{ij}}
= \langle u^*(e_j), e_i\rangle
= \langle e_j, u(e_i)\rangle = a_{ji} ,
\qquad\text{jadi}\qquad
B = \conj{A}^{\mathsf T} .
\] Adapun $u$ disebut
\emph{Hermitian}\index{endomorfisma Hermitian} apabila $u^* = u$
(yakni $A^\dagger = A$), \emph{uniter}\index{grup uniter} apabila $u^*u =
\mathrm{id}$ (yakni $A^\dagger A = I$: inilah grup $U(n)$), dan
\emph{normal} apabila $u^*u = uu^*$.
\end{definition}

\begin{proposition}\label{prop:b2:hermitian:eigenvalues}
\omterm{def:b2:reduction:eigen}{Nilai eigen} \omterm{def:b2:hermitian:adjoint}{endomorfisma Hermitian} bersifat \emph{real};
\omterm{def:b2:reduction:eigen}{nilai eigen} endomorfisma \omterm{def:b2:hermitian:adjoint}{uniter} \emph{bermodulus $1$}; dan pada
kedua kasusnya ruang eigen bagi \omterm{def:b2:reduction:eigen}{nilai eigen} yang berbeda saling ortogonal.
\end{proposition}

\begin{proof}
Kasus \omterm{def:b2:hermitian:adjoint}{Hermitian}, dengan $u(x) = \lambda x$ dan $x \neq 0$:
\[
\lambda \norm x^2 = \langle x, u(x)\rangle = \langle u(x), x\rangle
= \conj{\langle x, u(x)\rangle} = \conj\lambda\,\norm x^2 ,
\]
jadi $\lambda \in \R$. Kasus \omterm{def:b2:hermitian:adjoint}{uniter}: $\norm{u(x)} = \norm x$ (berkat $u^*u =
\mathrm{id}$), jadi $\abs\lambda\norm x = \norm x$. Keortogonalannya
(kasus \omterm{def:b2:hermitian:adjoint}{Hermitian}): $\lambda\langle x, y\rangle = \langle u(x),
y\rangle = \langle x, u(y)\rangle = \mu\langle x, y\rangle$ bagi
\omterm{def:b2:reduction:eigen}{vektor eigen} dengan $\lambda \neq \mu$ yang real. Adapun kasus \omterm{def:b2:hermitian:adjoint}{uniter},
selengkapnya: untuk $u(x) = \lambda x$ dan $u(y) = \mu y$ dengan $\lambda \neq
\mu$ (keduanya bermodulus satu),
\[
\langle x, y\rangle = \langle u(x), u(y)\rangle
= \conj\lambda\mu\,\langle x, y\rangle ,
\]
dan $\conj\lambda\mu = \frac{\mu}{\lambda} \neq 1$: jadi faktornya
bukan $1$, sehingga $\langle x, y\rangle = 0$.
\end{proof}

\begin{example}[Sebuah matriks anti-Hermitian, didiagonalkan]\label{ex:b2:hermitian:skewexample}
$A = \begin{pmatrix} 0 & -2\\ 2 & 0\end{pmatrix}$ memenuhi
$A^\dagger = A^{\mathsf T} = -A$: jadi \omterm{def:b2:hermitian:adjoint}{anti-Hermitian} (sekaligus
antisimetrik real --- dan atas $\R$ ia sama sekali tak punya \omterm{def:b2:reduction:eigen}{nilai eigen}).
\omterm{def:b2:reduction:charpoly}{Polinomial karakteristiknya} $X^2 + 4$: jadi \omterm{def:b2:reduction:eigen}{nilai eigennya} $\pm2\iu$,
yang imajiner murni, sebagaimana diramalkan \cref{exo:b2:hermitian:9} secara
umum. \omterm{def:b2:reduction:eigen}{Vektor eigennya}: $(A - 2\iu I)v = 0$ memberi $v_1 =
\frac{1}{\sqrt2}(1, \iu)$, lalu $v_2 = \frac{1}{\sqrt2}(1,
-\iu)$ untuk $-2\iu$; dan keduanya ortogonal:
\[
\langle v_1, v_2\rangle
= \tfrac12\bigl(\conj{1}\cdot1 +
\conj{\iu}\cdot(-\iu)\bigr)
= \tfrac12(1 - 1) = 0 .
\]
Jadi $A = U\operatorname{diag}(2\iu, -2\iu)\,U^\dagger$ dengan $U =
(v_1\ v_2)$ yang \omterm{def:b2:hermitian:adjoint}{uniter}. Pelajaran penutupnya: $H = -\iu A =
\begin{pmatrix} 0 & 2\iu\\ -2\iu & 0\end{pmatrix}$ bersifat \omterm{def:b2:hermitian:adjoint}{Hermitian}
dengan \omterm{def:b2:reduction:eigen}{spektrum} real $\{\pm2\}$ dan \omterm{def:b2:reduction:eigen}{vektor eigen} yang \emph{sama}
--- yakni bijeksi $u \mapsto \iu u$ antara \omterm{def:b2:hermitian:adjoint}{endomorfisma
Hermitian} dan \omterm{def:b2:hermitian:adjoint}{anti-Hermitian}
(\cref{exo:b2:hermitian:9}), yang terlihat matriks demi matriks; sedangkan atas $\R$
matriks $A$ yang sama adalah perputaran-penskalaan tanpa \omterm{def:b2:reduction:eigen}{vektor eigen} sama
sekali, dan hanya peralihan ke $\C$ yang menyingkap bentuk normalnya.
\end{example}

\begin{theorem}[Teorema spektral Hermitian]\label{thm:b2:hermitian:spectral}
Setiap \omterm{def:b2:hermitian:adjoint}{endomorfisma Hermitian} pada \omterm{def:b2:hermitian:def}{ruang Hermitian} punya basis ortonormal
beranggotakan \omterm{def:b2:reduction:eigen}{vektor eigen} (dengan \omterm{def:b2:reduction:eigen}{nilai eigen} yang real): jadi $A^\dagger = A$
mengakibatkan $A = U D U^{\dagger}$ dengan $U \in U(n)$ dan $D$
diagonal real.\index{teorema spektral!Hermitian}
\end{theorem}

\begin{proof}
Atas $\C$, \omterm{def:b2:reduction:charpoly}{polinomial karakteristiknya} terurai: jadi sebuah \omterm{def:b2:reduction:eigen}{vektor eigen}
$e_1$ ada (\cref{ch:b2:reduction}) --- tanpa perlu hujah kekompakan,
yakni salah satu keunggulan $\C$. Normalkan ia. Lalu komplemen
ortogonalnya $F = e_1^\perp$ bersifat stabil: sebab untuk $x \perp e_1$,
\[
\langle e_1, u(x)\rangle = \langle u(e_1), x\rangle
= \lambda_1\langle e_1, x\rangle = 0
\]
(dengan $\lambda_1$ real). Pembatasannya bersifat \omterm{def:b2:hermitian:adjoint}{Hermitian}; jadi
berinduksilah atas dimensinya lalu rangkaikan. Secara rinci: pembatasan $u|_F$
adalah endomorfisma \omterm{def:b2:hermitian:def}{ruang Hermitian} $F$ (yang berdimensi $n - 1$)
dengan $\langle u|_F(x), y\rangle = \langle x, u|_F(y)\rangle$
yang diwarisi dari $u$; lalu hipotesis induksinya menyediakan basis
ortonormal $(e_2, \dots, e_n)$ bagi $F$ yang beranggotakan \omterm{def:b2:reduction:eigen}{vektor eigen},
dan $(e_1, e_2, \dots, e_n)$ bersifat ortonormal di $E$ (sebab $e_1 \perp
F$) sekaligus tersusun atas \omterm{def:b2:reduction:eigen}{vektor eigen} $u$. Terjemahan matriksnya:
kolom $U$ adalah $e_i$, lalu $U^\dagger U = I$ menyatakan
keortonormalannya, dan $AU = UD$ mengumpulkan persamaan \omterm{def:b2:reduction:eigen}{nilai
eigennya}, sehingga $A = UDU^\dagger$ dengan $D$ diagonal real
(\cref{prop:b2:hermitian:eigenvalues}).
\end{proof}

\begin{example}\label{ex:b2:hermitian:pauli}
$A = \begin{pmatrix} 0 & -\iu\\ \iu & 0\end{pmatrix}$ bersifat \omterm{def:b2:hermitian:adjoint}{Hermitian}
(sebab $A^\dagger = A$): \omterm{def:b2:reduction:eigen}{nilai eigennya} dari $\chi_A = X^2 - 1$ adalah $\pm 1$
(yang real, seperti dijanjikan), dengan \omterm{def:b2:reduction:eigen}{vektor eigen} ortonormal
$\frac{1}{\sqrt2}(1, \iu)^{\mathsf T}$ dan $\frac{1}{\sqrt2}(1,
-\iu)^{\mathsf T}$. (Para fisikawan mengenal $A$ sebagai matriks Pauli; dan
kerealan \omterm{def:b2:reduction:eigen}{spektrum} Hermitianlah yang membuat besaran teramati kuantum dimodelkan
oleh operator \omterm{def:b2:hermitian:adjoint}{Hermitian}.)
\end{example}

\begin{example}[Sebuah matriks Hermitian definit positif, dikerjakan]\label{ex:b2:hermitian:pdrun}
$A = \begin{pmatrix} 2 & 1-\iu\\ 1+\iu & 3\end{pmatrix}$ bersifat
\omterm{def:b2:hermitian:adjoint}{Hermitian}, sebab diagonalnya real dan entri luar diagonalnya saling
sekawan. \omterm{def:b2:reduction:charpoly}{Polinomial karakteristiknya}:
\[
(2-\lambda)(3-\lambda) - \abs{1-\iu}^2
= \lambda^2 - 5\lambda + 4
= (\lambda - 1)(\lambda - 4) :
\]
jadi \omterm{def:b2:reduction:eigen}{spektrumnya} $\{1, 4\}$, yang real dan positif --- sehingga $A$ definit
positif. \omterm{def:b2:reduction:eigen}{Vektor eigennya}: untuk $\lambda = 4$, sistem $(A -
4I)v = 0$ memberi $v_4 = (1 - \iu,\ 2)$ (periksa baris keduanya:
$(1+\iu)(1-\iu) - 2 = 0$); sedangkan untuk $\lambda = 1$, $v_1 = (1 - \iu,\
-1)$. Keortogonalannya, dengan sekawannya pada slot pertama:
\[
\langle v_4, v_1\rangle
= \conj{(1-\iu)}\,(1-\iu) + \conj{2}\,(-1)
= 2 - 2 = 0 . \checkmark
\]
Setelah dinormalkan (dengan $\norm{v_4}^2 = 2 + 4 = 6$ dan $\norm{v_1}^2 = 2 + 1
= 3$) diperoleh $U = \bigl(\frac{v_4}{\sqrt6}\
\frac{v_1}{\sqrt3}\bigr)$ yang \omterm{def:b2:hermitian:adjoint}{uniter} dengan $A =
U\operatorname{diag}(4,1)U^\dagger$. Pelajaran penutupnya:
pembacaan Rayleighnya seketika --- sebab pada bola satuan
$\C^2$, besaran $\langle x, Ax\rangle$ berjangkau $\intcc{1}{4}$,
yang tercapai di kedua \omterm{def:b2:reduction:eigen}{vektor eigennya}; dan inilah benih $n = 2$ bagi
teori Courant--Fischer yang \omterm{def:b2:structures:generated}{dibangun} pada soal akhir pekan.
Periksa sekilas bahwa bentuknya real di luar \omterm{def:b2:reduction:eigen}{vektor eigennya} juga: di
$x = (1, \iu)$,
\[
Ax = \bigl(2 + (1-\iu)\iu,\ (1+\iu) + 3\iu\bigr)
= (3 + \iu,\ 1 + 4\iu),
\]
\[
\langle x, Ax\rangle = \conj{1}\,(3+\iu) +
\conj{\iu}\,(1+4\iu)
= (3 + \iu) + (-\iu)(1 + 4\iu) = 3 + \iu - \iu + 4 = 7 \in
\R ,
\]
sebagaimana dijamin mekanisme bukti \cref{prop:b2:hermitian:eigenvalues}
(yakni kesimetrikan-sekawan terhadap $A^\dagger = A$) bagi
setiap $x$.
\end{example}

\begin{example}[Sebuah matriks uniter didiagonalkan]\label{ex:b2:hermitian:unitary}
$U = \frac{1}{\sqrt2}\begin{pmatrix} 1 & \iu\\ \iu & 1
\end{pmatrix}$ (yang \omterm{def:b2:hermitian:adjoint}{uniter} menurut \cref{exo:b2:hermitian:2}). Adapun
\omterm{def:b2:reduction:charpoly}{polinomial karakteristiknya} $\bigl(X - \frac{1}{\sqrt2}\bigr)^2
+ \frac12$, dengan akar
\[
\lambda_\pm = \frac{1 \pm \iu}{\sqrt2} = \eu^{\pm\iu\pi/4},
\]
yang bermodulus satu sebagaimana dijanjikan \cref{prop:b2:hermitian:eigenvalues},
dan \omterm{def:b2:reduction:eigen}{vektor eigen} ortonormalnya $\frac{1}{\sqrt2}(1, \pm1)$. Jadi $U =
V\operatorname{diag}(\eu^{\iu\pi/4},
\eu^{-\iu\pi/4})V^\dagger$: dalam basis yang tepat, $U$ adalah sepasang
perputaran bidang sebesar $\pm\frac\pi4$ --- karena matriks perputaran real tak
punya \omterm{def:b2:reduction:eigen}{vektor eigen} real, dan atas $\C$ ia justru terurai menjadi dua
skalar bermodulus satu. Pelajaran penutupnya: \omterm{def:b2:reduction:eigen}{spektrum} \omterm{def:b2:hermitian:adjoint}{Hermitian} hidup pada
garis real, sedangkan \omterm{def:b2:reduction:eigen}{spektrum} \omterm{def:b2:hermitian:adjoint}{uniter} pada lingkaran satuan; keduanya
bayangan kenormalan yang sama, dan transformasi Cayley pada
\cref{exo:b2:hermitian:6} memetakan gambaran yang satu ke yang lain.
\end{example}

\begin{example}[Transformasi Cayley, dihitung]\label{ex:b2:hermitian:cayleyrun}
Jalankan \cref{exo:b2:hermitian:6} pada $H = \begin{pmatrix} 0 & 1\\
1 & 0\end{pmatrix}$ (yang \omterm{def:b2:hermitian:adjoint}{Hermitian}, berspektrum $\{1, -1\}$, dengan
\omterm{def:b2:reduction:eigen}{vektor eigen} ortonormal $\frac{1}{\sqrt2}(1, \pm1)$). Dalam basis
eigennya segalanya menjadi skalar: sebab transformasi $\lambda
\mapsto \frac{\lambda - \iu}{\lambda + \iu}$ mengirim
\[
1 \longmapsto \frac{1 - \iu}{1 + \iu} = -\iu,
\qquad
-1 \longmapsto \frac{-1 - \iu}{-1 + \iu} = \iu ,
\]
(kalikan dengan sekawan penyebutnya), jadi $U = (H -
\iu I)(H + \iu I)^{-1}$ adalah matriks \omterm{def:b2:hermitian:adjoint}{uniter} bernilai eigen
$\mp\iu$ pada \omterm{def:b2:reduction:eigen}{vektor eigen} yang sama itu:
\[
U = \frac{1}{2}\begin{pmatrix} 1 & 1\\ 1 & -1\end{pmatrix}
\begin{pmatrix} -\iu & 0\\ 0 & \iu\end{pmatrix}
\begin{pmatrix} 1 & 1\\ 1 & -1\end{pmatrix}
= \begin{pmatrix} 0 & -\iu\\ -\iu & 0\end{pmatrix} .
\]
Periksa: $U^\dagger U = I$, dan $1 \notin \operatorname{Sp}U =
\{\pm\iu\}$, sebagaimana dijanjikan teorinya. Pelajaran penutupnya: garis
realnya terpetakan ke lingkaran satuan dikurangi titik $1$ ---
\omterm{def:b2:reduction:eigen}{nilai eigen} demi \omterm{def:b2:reduction:eigen}{nilai eigen}, transformasi Cayley tak lain
pemetaan Möbius $\frac{\lambda-\iu}{\lambda+\iu}$, dan matriksnya
sekadar mengikuti \omterm{def:b2:reduction:eigen}{spektrumnya}.
\end{example}

\begin{remark}[Jebakan yang sering muncul]
\emph{(i) Tempat jatuhnya garis atas:} buku ini mengawankan slot
\emph{pertamanya}, jadi koordinatnya adalah $\langle e_i, x\rangle$
dan $\langle\lambda x, y\rangle = \conj\lambda\langle x,
y\rangle$; sedangkan banyak teks mengawankan slot keduanya ---
jadi terjemahkanlah dulu sebelum membandingkan rumus, atau tanda $\iu$
menjadi salah secara diam-diam. \emph{(ii) \omterm{def:b2:quadratic:def}{Polarisasi} kompleks lebih kuat:}
atas $\C$, bila $\langle x, u(x)\rangle = 0$ untuk \emph{setiap} $x$
maka $u = 0$ (uraikan $x + y$ dan $x + \iu y$: sebab bagian real
maupun bagian imajiner $\langle x, u(y)\rangle$ lenyap); sedangkan atas
$\R$ ini gagal --- sebab perputaran sebesar $\frac\pi2$ memenuhi
$\langle x, u(x)\rangle = 0$ di mana-mana. Karena itu, hanya atas
$\C$, syarat ``$\langle x, u(x)\rangle \in \R$ untuk setiap $x$''
sudah memaksa $u$ menjadi \omterm{def:b2:hermitian:adjoint}{Hermitian}. \emph{(iii) Normal real tidak
\omterm{def:b2:reduction:diag}{dapat didiagonalkan}:} matriks pada
\cref{ex:b2:hermitian:skewexample} bersifat normal tetapi tak punya \omterm{def:b2:reduction:eigen}{nilai
eigen} real; jadi pendiagonalan \omterm{def:b2:hermitian:adjoint}{uniter} adalah teorema \emph{atas}
$\C$, sedangkan atas $\R$ kita hanya memperoleh reduksi blok.
\emph{(iv) Memeriksa keuniteran:} syarat $U^\dagger U = I$ berarti
\emph{kolomnya} ortonormal terhadap hasil kali \omterm{def:b2:hermitian:adjoint}{Hermitiannya} ---
jadi menguji $UU^{\mathsf T}$, atau melupakan pengawanan pada
hasil kali kolomnya, adalah dua cara klasik untuk mengesahkan matriks
yang salah.
\end{remark}

\begin{example}[Isometri persis merupakan pemetaan uniter]\label{ex:b2:hermitian:isometry}
Pemeliharaan \omterm{def:b2:nvs:norm}{norma} tampak lebih lemah daripada keuniteran, tetapi atas $\C$
ia tidak: sebab bila $\norm{u(x)} = \norm x$ untuk setiap $x$, maka $u^*u =
\mathrm{id}$. Memang $v = u^*u - \mathrm{id}$ bersifat \omterm{def:b2:hermitian:adjoint}{Hermitian} dan
memenuhi $\langle x, v(x)\rangle = \norm{u(x)}^2 - \norm x^2 =
0$ untuk setiap $x$; jadi menurut \omterm{def:b2:quadratic:def}{polarisasi} kompleks (yakni jebakan (ii)
di atas), pemetaan yang ``diagonalnya'' lenyap secara identik pastilah nol:
sehingga $v = 0$. Secara konkret, \omterm{def:b2:quadratic:def}{polarisasinya} berbunyi
\[
0 = \langle x + y, v(x+y)\rangle
= \langle x, v(y)\rangle + \langle y, v(x)\rangle,
\qquad
0 = \langle x + \iu y, v(x + \iu y)\rangle
= \iu\langle x, v(y)\rangle - \iu\langle y, v(x)\rangle ,
\]
dan kedua barisnya bersama-sama memaksa $\langle x, v(y)\rangle = 0$
untuk setiap $x, y$. Pelajaran penutupnya: itulah sebabnya ``\omterm{def:b2:hermitian:adjoint}{uniter}'' dapat
diperiksa dengan mengukur panjangnya belaka --- yakni kekakuan yang kelak
dieksploitasi bab Fourier, tempat memelihara energi
$\norm f_2$ (yakni Parseval) sama dengan memelihara semua hasil kali
dalam koefisiennya.
\end{example}

\begin{remark}[Pandangan ke depan di dalam jilid ini]
Perkakas \omterm{def:b2:hermitian:adjoint}{Hermitian} yang \omterm{def:b2:structures:generated}{dibangun} di sini dikonsumsi hampir
seketika. Bab Fourier \emph{adalah} geometri \omterm{def:b2:hermitian:adjoint}{Hermitian}
dalam dimensi tak hingga: sebab eksponensial $(e_n)$ merupakan
keluarga ortonormal bagi $\langle f, g\rangle =
\frac{1}{2\pi}\int\conj fg$, ketaksamaan Bessel tak lain taksiran
proyeksi pada \cref{thm:b2:hermitian:cs} bab ini,
dan Parseval adalah kesamaan limitnya. Adapun transformasi Fourier hingga
(\cref{exo:b2:hermitian:10}) muncul kembali setiap kali konvolusi
harus didiagonalkan. Lalu soal akhir pekan bab ini ---
Courant--Fischer, Weyl, penyisipan --- memasok
kestabilan \omterm{def:b2:reduction:eigen}{nilai eigen} yang dipanggil bab persamaan diferensial
ketika ia menegaskan bahwa usikan kecil sebuah sistem
menggeser frekuensinya sedikit saja. Ke belakang, segalanya di sini
adalah cermin kompleks bab \omterm{def:b2:quadratic:def}{bentuk kuadratik}: jadi simpanlah
kedua kamusnya berdampingan ($A^{\mathsf T} \leftrightarrow
A^\dagger$, ortogonal $\leftrightarrow$ \omterm{def:b2:hermitian:adjoint}{uniter}, Rayleigh
real pada keduanya).
\end{remark}

\begin{remark}[Endomorfisma normal]
Atas $\C$ pernyataan yang definitif berbunyi: $u$ \emph{\omterm{def:b2:reduction:diag}{dapat
didiagonalkan} secara \omterm{def:b2:hermitian:adjoint}{uniter} bila dan hanya bila ia normal} (yakni $u^*u = uu^*$)
--- sehingga mencakup pemetaan \omterm{def:b2:hermitian:adjoint}{Hermitian}, \omterm{def:b2:hermitian:adjoint}{uniter} dan \omterm{def:b2:hermitian:adjoint}{anti-Hermitian} sekaligus. Adapun
buktinya merupakan penguatan yang menyenangkan atas hujah di atas
(\cref{exo:b2:hermitian:8}). Sebaliknya atas $\R$, kenormalan hanya
membeli pendiagonalan blok (yakni blok perputaran): jadi geometri kompleks
sungguh-sungguh lebih sederhana.
\end{remark}

\begin{remark}[Di mana ini dipakai]
Teori spektral \omterm{def:b2:hermitian:adjoint}{Hermitian} adalah matematika mekanika
kuantum: besaran teramati dimodelkan oleh operator \omterm{def:b2:hermitian:adjoint}{Hermitian}
(sebab \omterm{def:b2:reduction:eigen}{spektrum} real = nilai terukur), sedangkan evolusi waktunya oleh
operator \omterm{def:b2:hermitian:adjoint}{uniter} (sebab pemeliharaan \omterm{def:b2:nvs:norm}{norma} = kekekalan peluang). Di dalam
buku ini, bab Fourier bersandar pada keortonormalan
eksponensialnya --- yakni pernyataan tentang \omterm{def:b2:hermitian:def}{hasil kali dalam Hermitian} --- dan
pendiagonalan matriks sirkulan
(\cref{exo:b2:hermitian:10}) tak lain transformasi Fourier hingga.
Adapun soal akhir pekan mengembangkan kalkulus variasi bagi
\omterm{def:b2:reduction:eigen}{nilai eigen} (Courant--Fischer, Weyl, penyisipan), yakni roti
sehari-hari analisis numerik dan fisika matematis; sedangkan jilid Tahun
ke-3 memperluasnya ke operator \omterm{def:b2:hermitian:adjoint}{swa-adjoin} yang \omterm{def:b2:metric:compact}{kompak} pada
ruang Hilbert.
\end{remark}

\section{Latihan}

\begin{exercise}[$\star$]\label{exo:b2:hermitian:1}
Pada $\C^2$: hitunglah $\langle x, y\rangle$, $\norm x$, $\norm y$ untuk
$x = (1, \iu)$ dan $y = (\iu, 1)$; apakah keduanya ortogonal? Berikanlah
sebuah basis ortonormal yang memuat $\frac{x}{\norm x}$.
\end{exercise}

\begin{exercise}[$\star$]\label{exo:b2:hermitian:2}
Manakah yang \omterm{def:b2:hermitian:adjoint}{Hermitian}? yang \omterm{def:b2:hermitian:adjoint}{uniter}? yang normal?
\[
\begin{pmatrix} 1 & \iu\\ -\iu & 2 \end{pmatrix},
\qquad
\frac{1}{\sqrt2}\begin{pmatrix} 1 & \iu\\ \iu & 1\end{pmatrix},
\qquad
\begin{pmatrix} 0 & 1\\ 0 & 0 \end{pmatrix}.
\]
\end{exercise}

\begin{exercise}[$\star$]\label{exo:b2:hermitian:3}
Buktikan bahwa matriks $A \in \mathcal{M}_n(\C)$ dituliskan secara tunggal
$A = H + \iu K$ dengan $H, K$ \omterm{def:b2:hermitian:adjoint}{Hermitian} \emph{(yakni ``bagian real dan
bagian imajinernya'' $H = \frac{A + A^\dagger}{2}$, $K = \frac{A -
A^\dagger}{2\iu}$)}, dan bahwa $A$ normal bila dan hanya bila $H$ dan $K$
komut.
\end{exercise}

\begin{exercise}[$\star\star$]\label{exo:b2:hermitian:4}
Diagonalkan dalam basis ortonormal: $A = \begin{pmatrix} 2 & \iu\\
-\iu & 2\end{pmatrix}$, lalu hitunglah $A^k$ untuk $k \in \N$.
\end{exercise}

\begin{exercise}[$\star\star$]\label{exo:b2:hermitian:5}
Buktikan bahwa $U(n)$ \omterm{def:b2:metric:compact}{kompak}, dan bahwa pemetaan \omterm{def:b2:reduction:eigen}{nilai eigennya}
pada: setiap $\lambda$ yang bermodulus satu muncul bagi suatu matriks \omterm{def:b2:hermitian:adjoint}{uniter}.
Buktikan pula bahwa $\det U \in \mathbb{U}$ (yakni lingkaran satuan) untuk $U \in
U(n)$.
\end{exercise}

\begin{exercise}[$\star\star$]\label{exo:b2:hermitian:6}
(Transformasi Cayley) Misalkan $H$ \omterm{def:b2:hermitian:adjoint}{Hermitian}. Buktikan bahwa $H + \iu I$
terbalikkan dan bahwa $U = (H - \iu I)(H + \iu I)^{-1}$ bersifat \omterm{def:b2:hermitian:adjoint}{uniter},
dengan $1 \notin \operatorname{Sp}(U)$. \emph{(Bekerjalah secara spektral: sebab pada
basis eigen $H$, segalanya menjadi skalar.)}
\end{exercise}

\begin{exercise}[$\star\star$]\label{exo:b2:hermitian:7}
Untuk $A$ \omterm{def:b2:hermitian:adjoint}{Hermitian} definit positif (yakni $\langle x, Ax\rangle > 0$ untuk
$x \neq 0$), buktikan bahwa $\operatorname{Sp}(A) \subseteq
\intoo{0}{\infty}$, bahwa $A = B^2$ untuk suatu $B$ \omterm{def:b2:hermitian:adjoint}{Hermitian} definit
positif, dan bahwa $\det A > 0$.
\end{exercise}

\begin{exercise}[$\star\star\star$]\label{exo:b2:hermitian:8}
(Teorema spektral bagi endomorfisma normal) Misalkan $u$ normal pada sebuah
\omterm{def:b2:hermitian:def}{ruang Hermitian}.
\begin{enumerate}
  \item Buktikan $\norm{u(x)} = \norm{u^*(x)}$ untuk setiap $x$, lalu turunkan
        $\ker(u - \lambda) = \ker(u^* - \conj\lambda)$.
  \item Buktikan bahwa ruang eigen $u$ bagi \omterm{def:b2:reduction:eigen}{nilai eigen} yang berbeda saling
        ortogonal, dan bahwa komplemen ortogonal sebuah
        ruang eigen bersifat stabil terhadap $u$.
  \item Simpulkan secara induktif bahwa $u$ \omterm{def:b2:reduction:diag}{dapat didiagonalkan} secara \omterm{def:b2:hermitian:adjoint}{uniter};
        lalu sebaliknya.
\end{enumerate}
\end{exercise}

\begin{exercise}[$\star$]\label{exo:b2:hermitian:9}
Sebuah endomorfisma disebut \emph{\omterm{def:b2:hermitian:adjoint}{anti-Hermitian}} apabila $u^* = -u$. Buktikan
bahwa \omterm{def:b2:reduction:eigen}{nilai eigennya} imajiner murni, bahwa $u \mapsto \iu
u$ merupakan bijeksi dari \omterm{def:b2:hermitian:adjoint}{endomorfisma Hermitian} ke endomorfisma
\omterm{def:b2:hermitian:adjoint}{anti-Hermitian}, dan bahwa endomorfisma \omterm{def:b2:hermitian:adjoint}{anti-Hermitian} \omterm{def:b2:reduction:diag}{dapat
didiagonalkan} secara \omterm{def:b2:hermitian:adjoint}{uniter} (\cref{exo:b2:hermitian:8}).
\end{exercise}

\begin{exercise}[$\star\star$]\label{exo:b2:hermitian:10}
(Transformasi Fourier hingga) Misalkan $S$ geseran \omterm{def:b2:structures:generated}{siklik} pada
$\C^n$: $S(x_0, x_1, \dots, x_{n-1}) = (x_{n-1}, x_0, \dots,
x_{n-2})$, dan $\omega = \eu^{2\iu\pi/n}$.
\begin{enumerate}
  \item Tunjukkan bahwa $S$ \omterm{def:b2:hermitian:adjoint}{uniter}, dan bahwa vektor $f_k =
        \frac{1}{\sqrt n}\bigl(1, \omega^k, \omega^{2k}, \dots,
        \omega^{(n-1)k}\bigr)$ dengan $0 \leq k < n$ membentuk
        basis ortonormal beranggotakan \omterm{def:b2:reduction:eigen}{vektor eigen}: $Sf_k = \omega^{-k}
        f_k$.
  \item Turunkan bahwa setiap matriks \emph{sirkulan} $C =
        \sum_{j=0}^{n-1} c_jS^j$ bersifat normal, didiagonalkan oleh
        basis yang sama, dengan \omterm{def:b2:reduction:eigen}{nilai eigen} $\widehat c(k) = \sum_j
        c_j\,\omega^{-jk}$.
\end{enumerate}
\end{exercise}

\begin{exercise}[$\star\star$]\label{exo:b2:hermitian:11}
Misalkan $P$ endomorfisma idempoten (yakni $P^2 = P$) pada sebuah \omterm{def:b2:hermitian:def}{ruang
Hermitian}. Buktikan bahwa $P$ adalah proyeksi \emph{ortogonal} ke
$\operatorname{im} P$ bila dan hanya bila $P^* = P$. Berikanlah matriks
proyeksi ortogonal ke $\C v$ (dengan $\norm v = 1$), dan
ke sebuah subruang berbasis ortonormal $(v_1, \dots, v_k)$.
\end{exercise}

\begin{exercise}[$\star\star\star$]\label{exo:b2:hermitian:12}
(Proyektor spektral lewat interpolasi) Misalkan $A$ \omterm{def:b2:hermitian:adjoint}{Hermitian} dengan
\omterm{def:b2:reduction:eigen}{nilai eigen} yang \emph{berbeda-beda} $\lambda_1, \dots, \lambda_p$ dan
penguraian ruang eigen $E = \bigoplus_i E_i$. Definisikan
polinomial Lagrange $L_i(X) = \prod_{j\neq i}\frac{X -
\lambda_j}{\lambda_i - \lambda_j}$. Buktikan bahwa $P_i = L_i(A)$
adalah proyeksi ortogonal ke $E_i$, bahwa $P_iP_j = 0$ untuk
$i \neq j$, bahwa $\sum_i P_i = I$, dan bahwa $A = \sum_i \lambda_iP_i$
(yakni \emph{dekomposisi spektralnya}); lalu nyatakanlah $f(A)$ untuk sembarang
polinomial $f$ dalam $P_i$.
\end{exercise}

\section{Soal: Courant--Fischer, Weyl, dan kalkulus
nilai eigen}

\begin{problem}\label{pb:b2:hermitian:1}
\omterm{def:b2:reduction:eigen}{Nilai eigen} sebuah matriks \omterm{def:b2:hermitian:adjoint}{Hermitian} bukan sekadar akar sebuah
polinomial: ia penyelesaian \emph{masalah pengoptimalan}.
Sudut pandang variasional itu --- yakni \omterm{pb:b2:hermitian:1}{hasil bagi Rayleigh} dan
\emph{teorema min-maks Courant--Fischer}\index{teorema
Courant--Fischer} --- membuat \omterm{def:b2:reduction:eigen}{nilai eigen} dapat dibandingkan, stabil, dan
terhitung, dan soal ini memanen hasil klasiknya:
\emph{ketaksamaan usikan Weyl}\index{ketaksamaan
Weyl}, \emph{penyisipan Cauchy}, ketaksamaan trace Schur dan Ky Fan,
kemonotonan akar kuadrat matriks,
dan \omterm{def:b2:reduction:eigen}{spektrum} Laplacian diskret. Di sepanjang soal ini, $A, B,
E$ bersifat \omterm{def:b2:hermitian:adjoint}{Hermitian} pada $E = \C^n$ dengan \omterm{def:b2:reduction:eigen}{nilai eigen} yang didaftar
menurun $\lambda_1(A) \geq \dots \geq \lambda_n(A)$,
dan $R_A(x) = \frac{\langle x, Ax\rangle}{\langle x, x\rangle}$
untuk $x \neq 0$ disebut \emph{hasil bagi Rayleigh}\index{hasil bagi
Rayleigh}.

\medskip
\noindent\textbf{Bagian I --- \omterm{pb:b2:hermitian:1}{Hasil bagi Rayleigh} dan min-maks.}
Tetapkan sebuah basis eigen ortonormal $(e_1, \dots, e_n)$ dengan $Ae_i =
\lambda_ie_i$.
\begin{enumerate}
  \item Tunjukkan bahwa $R_A(x)$ bersifat real, dan bahwa
        \[
        \lambda_n \leq R_A(x) \leq \lambda_1
        \qquad (x \neq 0),
        \]
        dengan kedua batasnya tercapai: $\lambda_1 = \max R_A$ dan
        $\lambda_n = \min R_A$.
  \item Tunjukkan bahwa titik kritis $R_A$ persis merupakan
        \omterm{def:b2:reduction:eigen}{vektor eigen} $A$ \emph{(uraikan $t \mapsto R_A(x +
        tv)$ di $t = 0$ untuk $v$ sembarang, lalu gantilah $v$ dengan
        $\iu v$)}.
  \item Misalkan $V_k = \operatorname{Vect}(e_1, \dots, e_k)$ dan
        $W_k = \operatorname{Vect}(e_k, \dots, e_n)$. Tunjukkan
        \[
        \min_{x \in V_k\setminus\{0\}} R_A(x) = \lambda_k
        = \max_{x \in W_k\setminus\{0\}} R_A(x) .
        \]
  \item Buktikan \emph{teorema Courant--Fischer}: untuk $1 \leq
        k \leq n$,
        \[
        \lambda_k = \max_{\dim V = k}\;\min_{x \in
        V\setminus\{0\}} R_A(x)
        = \min_{\dim W = n-k+1}\;\max_{x \in
        W\setminus\{0\}} R_A(x)
        \]
        \emph{(sebab untuk sembarang $V$ berdimensi $k$ berlaku $V \cap W_k \neq
        \{0\}$ menurut Grassmann, jadi $\min_V R_A \leq \lambda_k$;
        lalu pertanyaan 3 menunjukkan batasnya tercapai)}.
  \item (Kemonotonan) Tulislah $A \leq B$ apabila $B - A$ bersifat
        semidefinit positif. Turunkan dari pertanyaan 4: bahwa $A \leq
        B$ mengakibatkan $\lambda_k(A) \leq \lambda_k(B)$ untuk setiap
        $k$.
\end{enumerate}

\medskip
\noindent\textbf{Bagian II --- Ketaksamaan Weyl.}
\begin{enumerate}[resume]
  \item Tunjukkan bahwa subruang $V, W \subseteq \C^n$ dengan $\dim V
        + \dim W > n$ berpotongan secara tak trivial, lalu samaratakan:
        $\dim(V_1 \cap V_2 \cap V_3) \geq \dim V_1 + \dim V_2 +
        \dim V_3 - 2n$.
  \item Buktikan \emph{ketaksamaan Weyl}: untuk $i + j - 1 \leq n$,
        \[
        \lambda_{i+j-1}(A + B) \leq \lambda_i(A) +
        \lambda_j(B)
        \]
        \emph{(potongkanlah subruang $W_i(A)$, $W_j(B)$ dan
        $V_{i+j-1}(A+B)$ pada pertanyaan 3 lalu cacahlah
        dimensinya)}.
  \item Definisikan $\vertiii{E}_2 = \max_{\norm x = 1}\norm{Ex}$
        lalu tunjukkan $\vertiii E_2 = \max_k\abs{\lambda_k(E)}$ untuk
        $E$ yang \omterm{def:b2:hermitian:adjoint}{Hermitian}. Turunkan \emph{teorema usikan
        Weyl}:
        \[
        \bigl|\lambda_k(A + E) - \lambda_k(A)\bigr| \leq
        \vertiii{E}_2
        \qquad (1 \leq k \leq n) :
        \]
        jadi tiap \omterm{def:b2:reduction:eigen}{nilai eigennya} merupakan fungsi \omterm{def:b2:metric:continuity}{Lipschitz} berkonstanta
        $1$ terhadap matriksnya.
  \item (Usikan berrank satu) Misalkan $P$ \omterm{def:b2:hermitian:adjoint}{Hermitian} semidefinit
        positif berrank $1$. Tunjukkan
        \[
        \lambda_k(A) \leq \lambda_k(A + P) \leq
        \lambda_{k-1}(A) \qquad (2 \leq k \leq n),
        \]
        bersama $\lambda_1(A) \leq \lambda_1(A+P)$: jadi
        \omterm{def:b2:reduction:eigen}{nilai eigen} barunya \emph{menyisip} di antara yang lama.
  \item Periksalah pertanyaan 8 secara numerik: $A = \begin{pmatrix} 2 &
        \iu\\ -\iu & 2\end{pmatrix}$
        (\cref{exo:b2:hermitian:4}: berspektrum $\{3, 1\}$) dan
        $E = \begin{pmatrix} 0 & 1\\ 1 & 0\end{pmatrix}$
        (berspektrum $\{1, -1\}$): hitunglah \omterm{def:b2:reduction:eigen}{spektrum} $A +
        E$ beserta kedua ruas ketaksamaannya.
\end{enumerate}

\medskip
\noindent\textbf{Bagian III --- Penyisipan dan ketaksamaan
trace.}
\begin{enumerate}[resume]
  \item (Penyisipan Cauchy) Misalkan $B$ submatriks utama terdepan
        berukuran $(n-1)\times(n-1)$ milik $A$. Buktikan
        \[
        \lambda_{k+1}(A) \leq \lambda_k(B) \leq \lambda_k(A)
        \qquad (1 \leq k \leq n-1)
        \]
        \emph{(pandanglah $\C^{n-1} \subseteq \C^n$; sebab di atasnya, $R_B$
        adalah pembatasan $R_A$; lalu terapkan Courant--Fischer pada
        kedua arasnya)}.
  \item Iterasikan: untuk submatriks utama $B$ berukuran $n - m$,
        berlaku $\lambda_{k+m}(A) \leq \lambda_k(B) \leq
        \lambda_k(A)$.
  \item (Schur) Misalkan $d_1 \geq d_2 \geq \dots \geq d_n$ menyatakan
        entri diagonal $A$ yang telah disortir. Buktikan, untuk setiap $k$:
        \[
        \sum_{i=1}^{k} d_i \leq \sum_{i=1}^{k}\lambda_i(A),
        \]
        dengan kesamaan di $k = n$ (yakni tracenya)
        \emph{(sebab $k$ entri diagonal terpilihnya membentuk
        submatriks utama $k\times k$; lalu batasilah tracenya lewat
        pertanyaan 12)}.
  \item (Ky Fan) Buktikan
        \[
        \sum_{i=1}^{k}\lambda_i(A)
        = \max\Bigl\{\sum_{i=1}^{k}\langle x_i, Ax_i\rangle
        : (x_1, \dots, x_k) \text{ ortonormal}\Bigr\} .
        \]
  \item Periksalah pertanyaan 11 dan 13 pada
        \[
        A = \begin{pmatrix} 2 & 1 & 0\\ 1 & 2 & 1\\
        0 & 1 & 2\end{pmatrix}
        \qquad
        \bigl(\operatorname{Sp} = \{2 - \sqrt2,\ 2,\
        2+\sqrt2\}\bigr)
        \]
        terhadap blok terdepan $2\times2$-nya (yang berspektrum $\{1,
        3\}$) dan terhadap diagonalnya.
\end{enumerate}

\medskip
\noindent\textbf{Bagian IV --- Urutan Loewner.} Lambang $A \leq B$
tetap berarti $B - A$ semidefinit positif; dan semua matriks pada
bagian ini bersifat \omterm{def:b2:hermitian:adjoint}{Hermitian}.
\begin{enumerate}[resume]
  \item Tunjukkan: bahwa $A \leq B$ mengakibatkan $a_{ii} \leq b_{ii}$ untuk setiap
        $i$, bahwa $\operatorname{tr} A \leq \operatorname{tr} B$,
        dan bahwa $C^\dagger AC \leq C^\dagger BC$ untuk \emph{setiap}
        matriks kompleks $C$.
  \item Tunjukkan bahwa pengkuadratan \emph{tidak} monoton: sebab untuk
        \[
        A = \begin{pmatrix} 1 & 0\\ 0 & 0\end{pmatrix},
        \qquad
        B = \begin{pmatrix} 2 & 1\\ 1 & 1\end{pmatrix},
        \]
        periksalah $0 \leq A \leq B$ tetapi $A^2 \not\leq B^2$.
  \item Buktikan bahwa akar kuadratnya \emph{memang} monoton: bahwa $0 \leq
        A \leq B$ mengakibatkan $\sqrt A \leq \sqrt B$ \emph{(misalkan
        $\mu$ \omterm{def:b2:reduction:eigen}{nilai eigen} $\sqrt B - \sqrt A$ dengan \omterm{def:b2:reduction:eigen}{vektor eigen}
        satuan $v$; lalu hitunglah $\langle v, (B - A)v\rangle =
        \mu\bigl(\langle v, \sqrt B\,v\rangle + \langle v,
        \sqrt A\,v\rangle\bigr)$ lalu bahaslah)}.
  \item Buktikan bahwa pembalikannya bersifat antitonik pada matriks
        definit positif: bahwa $0 < A \leq B$ mengakibatkan $B^{-1} \leq A^{-1}$
        \emph{(kongruenkan dengan $A^{-1/2}$ untuk mereduksinya ke $I \leq M
        \Rightarrow M^{-1} \leq I$, yang menjadi skalar dalam basis
        spektralnya)}.
  \item Misalkan $A, B$ definit positif. Tunjukkan bahwa
        \omterm{def:b2:reduction:eigen}{nilai eigen} $AB$ (yang secara umum tak \omterm{def:b2:hermitian:adjoint}{Hermitian}!) bersifat
        real dan positif, dan bahwa
        \[
        \lambda_{\max}(AB) \leq
        \lambda_{\max}(A)\,\lambda_{\max}(B)
        \]
        \emph{(konjugasikan dengan $\sqrt A$: sebab $AB \sim \sqrt
        A\,B\sqrt A$)}.
\end{enumerate}

\medskip
\noindent\textbf{Bagian V --- Laplacian diskret, dikerjakan.}
Misalkan $T_n$ matriks tridiagonal $n \times n$ dengan $2$ pada
diagonalnya dan $-1$ pada kedua diagonal tetangganya.
\begin{enumerate}[resume]
  \item Dengan $\theta_k = \frac{k\pi}{n+1}$, periksalah bahwa
        vektor $v_k = \bigl(\sin(j\theta_k)\bigr)_{1\leq j\leq
        n}$ memenuhi $T_nv_k = (2 - 2\cos\theta_k)\,v_k$
        \emph{(lewat kesamaan hasil-kali-ke-jumlah; lalu periksa baris
        perbatasannya $j = 1, n$)}. Simpulkan:
        \[
        \operatorname{Sp}(T_n) = \Bigl\{4\sin^2
        \frac{k\pi}{2(n+1)} : 1 \leq k \leq n\Bigr\},
        \]
        yang semuanya sederhana dan positif: jadi $T_n$ definit positif.
  \item (Sebuah potensial) Untuk $D =
        \operatorname{diag}(d_1, \dots, d_n)$ diagonal real, apitlah
        \omterm{def:b2:reduction:eigen}{spektrumnya}: untuk setiap $k$,
        \[
        \lambda_k(T_n) + \min_i d_i \;\leq\;
        \lambda_k(T_n + D) \;\leq\; \lambda_k(T_n) + \max_i
        d_i .
        \]
  \item Periksalah penyisipan Cauchy antara $T_3$ dan $T_2$
        secara gamblang (dengan \omterm{def:b2:reduction:eigen}{spektrum} $\{2 \pm \sqrt2, 2\}$ dan $\{1,
        3\}$), lalu tafsirkanlah: sebab $T_2$ adalah $T_3$ dengan satu ujung
        lintasannya dibuang.
  \item Tunjukkan bahwa \omterm{def:b2:reduction:eigen}{nilai eigen} ekstremnya memenuhi, ketika $n \to
        \infty$:
        \[
        \lambda_{\min}(T_n) = 4\sin^2\frac{\pi}{2(n+1)}
        \sim \frac{\pi^2}{(n+1)^2},
        \qquad
        \lambda_{\max}(T_n) \to 4 ,
        \]
        jadi bilangan kondisinya $\kappa_n =
        \lambda_{\max}/\lambda_{\min}$ tumbuh seperti
        $\frac{4(n+1)^2}{\pi^2}$: sehingga mendiskretkan turunan
        kedua pada kisi yang makin halus bersifat
        buruk-kondisi secara hakiki.
  \item Rangkuman. Satu kalimat untuk masing-masing: (i) mengapa
        pencirian variasionalnya, bukan \omterm{def:b2:reduction:charpoly}{polinomial
        karakteristiknya}, yang membuat \omterm{def:b2:reduction:eigen}{nilai eigen} stabil
        (pertanyaan 8--9); (ii) pertanyaan mana yang hanya memakai
        $\lambda_1 = \max R_A$ dan mana yang menuntut min-maks
        selengkapnya; (iii) apa yang ditambahkan urutan Loewner pada
        ceritanya; (iv) di mana perkakas ini muncul kembali (yakni pada
        analisis numerik matriks kekakuan pertanyaan 24; pada teori
        usikan kuantum; dan, pada jilid Tahun ke-3, pada asas
        min-maks bagi operator \omterm{def:b2:hermitian:adjoint}{swa-adjoin} yang
        \omterm{def:b2:metric:compact}{kompak}).

\end{enumerate}
\end{problem}
